% =========================================================
% Complex Analysis — Problems + Background + Solutions
% (branch-cut contour, Fourier transforms, SC rectangle map)
% Step-by-step style; no custom environments; prints cleanly
% =========================================================
\documentclass[11pt]{article}

\usepackage[a4paper,margin=1in]{geometry}
\usepackage{amsmath,amssymb,mathtools}
\usepackage{microtype}
\setlength{\parskip}{0.85em}
\setlength{\parindent}{0pt}

\DeclareMathOperator{\Res}{Res}

\newcommand{\BoxEq}[1]{\begin{equation*}\boxed{\displaystyle #1}\end{equation*}}

\DeclareMathOperator{\arccot}{arccot}
\DeclareMathOperator{\arsinh}{arsinh}
\DeclareMathOperator{\arcosh}{arcosh}
\DeclareMathOperator{\artanh}{artanh}
\DeclareMathOperator{\arcoth}{arcoth}

\usepackage{tikz,pgfplots}
\pgfplotsset{compat=1.18}

\usepackage{enumitem}

\begin{document}
	
	\begin{center}
		{\Large\bf Complex Analysis: Contours, Fourier Transforms, and Schwarz--Christoffel}\\[0.45em]
		{\small Problems with background and complete solutions}
	\end{center}
	
	\bigskip
	
	% =========================================================
	\section{Problem 1: Evaluate \(\displaystyle \int_{-1}^{1}\frac{\sqrt{1-x^{2}}}{x^{2}+1}\,dx\) by contour integration}
	
	\textit{Hint.} Consider \(F(z)=\sqrt{z^{2}-1}/(z^{2}+1)\) with a branch cut on \([-1,1]\), and integrate around a large circle.
	
	\subsection*{Background}
	Fix the branch of \(G(z)=\sqrt{z^{2}-1}\) with cut on \([-1,1]\) and \(G(x)>0\) for \(x>1\).
	For \(x\in(-1,1)\),
	\[
	G(x+i0)=+\,i\sqrt{1-x^{2}},\qquad G(x-i0)=-\,i\sqrt{1-x^{2}}.
	\]
	Hence the sum of the upper and lower lip integrals across the cut equals
	\[
	2i\int_{-1}^{1}\frac{\sqrt{1-x^{2}}}{x^{2}+1}\,dx = 2i\,I.
	\]
	
	\subsection*{Solution}
	Define
	\[
	F(z)=\frac{\sqrt{z^{2}-1}}{z^{2}+1},\qquad z\in\mathbb C\setminus[-1,1].
	\]
	Poles: \(z=\pm i\) (simple). Large-circle behavior: \(F(z)=\dfrac{1}{z}+O(z^{-3})\); thus
	\[
	\int_{|z|=R}F(z)\,dz \longrightarrow 2\pi i \qquad (R\to\infty).
	\]
	Deforming the big circle to a contour hugging the cut gives
	\[
	\underbrace{\int_{|z|=R}F(z)\,dz}_{\to\,2\pi i}
	=
	\underbrace{\int_{\text{upper lip}}F\,dz+\int_{\text{lower lip}}F\,dz}_{=\,2i\,I}
	+\text{(vanishing end caps)}.
	\]
	Residues at the simple poles:
	\begin{align*}
		\Res_{z=i}F
		&=\lim_{z\to i}\frac{(z-i)\sqrt{z^{2}-1}}{(z-i)(z+i)}
		=\frac{\sqrt{i^{2}-1}}{2i}
		=\frac{\sqrt{-2}}{2i}
		=\frac{\sqrt2}{2},\\
		\Res_{z=-i}F
		&=\lim_{z\to -i}\frac{(z+i)\sqrt{z^{2}-1}}{(z-i)(z+i)}
		=\frac{\sqrt{(-i)^{2}-1}}{-2i}
		=\frac{\sqrt{-2}}{-2i}
		=\frac{\sqrt2}{2}.
	\end{align*}
	Hence, by the residue theorem,
	\[
	\int_{|z|=R}F(z)\,dz \longrightarrow 2\pi i\left(\frac{\sqrt2}{2}+\frac{\sqrt2}{2}\right)=2\pi i\sqrt2.
	\]
	Equating the two expressions for the big-circle integral,
	\[
	2i\,I=2\pi i(\sqrt2-1)\quad\Rightarrow\quad
	\boxed{\,I=\pi(\sqrt2-1)\, }.
	\]
	
	\subsection*{Check (real calculus)}
	With \(x=\sin\theta\),
	\[
	I=2\int_{0}^{\pi/2}\frac{\cos^{2}\theta}{1+\sin^{2}\theta}\,d\theta
	=2\int_{0}^{\infty}\frac{dt}{(1+t^{2})(1+2t^{2})}\,dt
	=\pi(\sqrt2-1).
	\]
	
	\bigskip\hrule\bigskip
	
	% =========================================================
	\section{Problem 2: Fourier transforms and inverse verification}
	
	We use the convention \(\displaystyle \widehat f(\xi)=\int_{-\infty}^{\infty} f(x)\,e^{-i\xi x}\,dx\) and
	\(\displaystyle f(x)=\frac{1}{2\pi}\int_{-\infty}^{\infty}\widehat f(\xi)\,e^{i\xi x}\,d\xi\).
	
	\subsection*{(a) \(f(x)=e^{-|x|}\) — Transform and step-by-step inverse}
	\textit{Transform:}
	\[
	\widehat f(\xi)
	=\int_{0}^{\infty}e^{-(1+i\xi)x}\,dx+\int_{-\infty}^{0}e^{(1-i\xi)x}\,dx
	=\frac{1}{1+i\xi}+\frac{1}{1-i\xi}
	=\boxed{\frac{2}{1+\xi^{2}}}.
	\]
	
	\textit{Inverse by residues (two cases):}
	\begin{align*}
		f(x)
		&=\frac{1}{2\pi}\int_{\mathbb R}\frac{2\,e^{i\xi x}}{1+\xi^{2}}\,d\xi
		=\frac{1}{2\pi}\int_{\mathbb R}\frac{2\,e^{i\xi x}}{(\xi-i)(\xi+i)}\,d\xi.
		\\[0.25em]
		\textbf{Case }x>0:\quad
		f(x)
		&=\frac{1}{2\pi}\oint_{\Gamma^+}\frac{2\,e^{i\xi x}}{(\xi-i)(\xi+i)}\,d\xi
		&&\text{(close upward; Jordan’s lemma)}\\
		&=\frac{1}{2\pi}(2\pi i)\,\Res_{\xi=i}\!\left(\frac{2\,e^{i\xi x}}{(\xi-i)(\xi+i)}\right)
		=\frac{i\,2\,e^{ix}}{2i( i+i)}=\;e^{-x}.
		\\[0.35em]
		\textbf{Case }x<0:\quad
		f(x)
		&=\frac{1}{2\pi}\oint_{\Gamma^-}\frac{2\,e^{i\xi x}}{(\xi-i)(\xi+i)}\,d\xi
		&&\text{(close downward)}\\
		&=\frac{1}{2\pi}(2\pi i)\,\Res_{\xi=-i}\!\left(\frac{2\,e^{i\xi x}}{(\xi-i)(\xi+i)}\right)
		=\;e^{\,x}.
	\end{align*}
	Therefore \(\boxed{\,f(x)=e^{-|x|}\,}\).
	
	\subsection*{(b) \(f(x)=e^{-a^{2}x^{2}}\) (\(a>0\)) — Transform and step-by-step inverse}
	\textit{Transform:}
	\begin{align*}
		\widehat f(\xi)
		&=\int_{\mathbb R}e^{-a^{2}x^{2}-i\xi x}\,dx
		= e^{-\xi^{2}/(4a^{2})}\int_{\mathbb R}e^{-a^{2}\left(x+i\frac{\xi}{2a^{2}}\right)^{2}}dx\\
		&\phantom{=} \qquad \Rightarrow\qquad
		\boxed{\,\widehat f(\xi)=\frac{\sqrt{\pi}}{a}\,e^{-\xi^{2}/(4a^{2})}\, }.
	\end{align*}
	(Justify by integrating over a large rectangle whose vertical sides are the real line and
	the line shifted up by \(i\,\xi/(2a^{2})\); the horizontal sides vanish by Gaussian decay.)
	
	\textit{Inverse by a rectangular shift (explicit steps):}
	\begin{align*}
		f(x)
		&=\frac{1}{2\pi}\int_{\mathbb R}\frac{\sqrt{\pi}}{a}\,e^{-\xi^{2}/(4a^{2})}\,e^{i\xi x}\,d\xi\\
		&=\frac{1}{2\pi}\cdot\frac{\sqrt{\pi}}{a}\,
		\int_{\mathbb R} \exp\!\left(-\frac{1}{4a^{2}}(\xi-2ia^{2}x)^{2}\right)\,
		\exp\!\left(-a^{2}x^{2}\right)\,d\xi
		&&\text{(complete the square)}\\
		&=\frac{1}{2\pi}\cdot\frac{\sqrt{\pi}}{a}\,e^{-a^{2}x^{2}}
		\int_{\mathbb R+2ia^{2}x} \exp\!\left(-\frac{u^{2}}{4a^{2}}\right)\,du
		&&\text{(shift contour vertically)}\\
		&=\frac{1}{2\pi}\cdot\frac{\sqrt{\pi}}{a}\,e^{-a^{2}x^{2}}
		\int_{\mathbb R} \exp\!\left(-\frac{u^{2}}{4a^{2}}\right)\,du
		&&\text{(no poles; horizontal sides vanish)}\\
		&=\frac{1}{2\pi}\cdot\frac{\sqrt{\pi}}{a}\,e^{-a^{2}x^{2}}\cdot 2a\sqrt{\pi}
		\;=\;e^{-a^{2}x^{2}}.
	\end{align*}
	
	\bigskip\hrule\bigskip
	
	% =========================================================
	\section{Problem 3: Schwarz--Christoffel map from the half-plane to a rectangle}
	
	\subsection*{Statement}
	Write down, as an integral, the Schwarz--Christoffel map from the upper half-plane \(\mathbb H\) to a rectangle,
	with vertices being the images of the points \(z=\pm1\) and \(z=\pm a\) (with \(a>1\) real).
	Explain why \(a\) cannot be chosen arbitrarily, but is determined by the rectangle’s aspect ratio.
	
	\subsection*{Background}
	For prevertices \(x_k\in\mathbb R\) mapping to interior angles \(\alpha_k\pi\) in the polygon,
	\[
	f'(z)=C\prod_k (z-x_k)^{\alpha_k-1}\qquad\text{(up to an affine change of variables).}
	\]
	A rectangle has \(\alpha_k=\tfrac12\) at each vertex, so every factor’s exponent is \(-\tfrac12\).
	
	\subsection*{Solution}
	Place prevertices at \(-a,-1,1,a\) (in order) with \(a>1\).
	Then
	\[
	f'(z)=C\,(z{+}a)^{-1/2}(z{+}1)^{-1/2}(z{-}1)^{-1/2}(z{-}a)^{-1/2}
	=\frac{C}{\sqrt{(z^{2}-1)(z^{2}-a^{2})}}.
	\]
	Integrating,
	\[
	\boxed{\qquad f(z)=A+C\int^{\,z}\frac{dt}{\sqrt{(t^{2}-1)(t^{2}-a^{2})}}\qquad}
	\]
	maps \(\mathbb H\) conformally to a rectangle (up to translation and scaling).
	
	\subsection*{Aspect ratio determines \(a\)}
	Side lengths are proportional to real integrals between consecutive prevertices:
	\[
	W(a)=\int_{-1}^{1}\frac{dt}{\sqrt{(1-t^{2})(a^{2}-t^{2})}},\qquad
	H(a)=\int_{1}^{a}\frac{dt}{\sqrt{(t^{2}-1)(a^{2}-t^{2})}}.
	\]
	Let \(k=1/a\in(0,1)\). Standard substitutions give
	\[
	W(a)=\frac{2}{a}\,K(k),\qquad H(a)=\frac{2}{a}\,K'(k),
	\]
	where \(K\) is the complete elliptic integral of the first kind and \(K'(k)=K(\sqrt{1-k^{2}})\).
	Thus
	\[
	\frac{H(a)}{W(a)}=\frac{K'(k)}{K(k)}\qquad\text{with }k=\frac{1}{a}.
	\]
	Given a rectangle, this ratio fixes \(k\) (hence \(a\)) uniquely; therefore \(a\) is not free.
	
	\bigskip
	
	
	% ==================== Part (a) ====================
	\section*{4(a). Statement and branch points}
	Suppose \(\zeta^2=z^2+1\) and
	\[
	z-i=r_1e^{i\theta_1},\qquad z+i=r_2e^{i\theta_2}\qquad (r_1,r_2>0).
	\]
	Then
	\[
	\boxed{\ \zeta=\pm(r_1r_2)^{1/2}e^{\,i(\theta_1+\theta_2)/2}\ }.
	\]
	\emph{Reason.} \(z^2+1=(z-i)(z+i)=r_1r_2e^{i(\theta_1+\theta_2)}\). Any square root has the two signs above.
	The branch points of the multifunction \((z^2+1)^{1/2}\) are where either factor vanishes, i.e. \(z=\pm i\).
	
	% ==================== Part (b) ====================
	\section*{4(b). A branch with a finite vertical cut}
	Define the branch by
	\[
	f(z)=(r_1r_2)^{1/2}e^{\,i(\theta_1+\theta_2)/2},\qquad -\frac{\pi}{2}<\theta_1,\theta_2\le\frac{3\pi}{2}.
	\]
	
	\subsection*{Values across the imaginary axis}
	Let \(z=x+iy\) and set \(\phi(y)=\tfrac12(\theta_1+\theta_2)\).
	Note
	\[
	r_1=|y-1|,\quad r_2=|y+1|,\qquad
	\arg(i t)=\begin{cases}\pi/2,&t>0,\\ -\pi/2,&t<0.\end{cases}
	\]
	\textbf{(i) \(|y|<1\).} \(y-1<0\), \(y+1>0\).
	\begin{align*}
		x\to 0^+:&\ \theta_1\to(-\pi/2)^+,\ \theta_2\to(\pi/2)^- \Rightarrow \phi(y)\to 0,\\
		x\to 0^-:&\ \theta_1\to(3\pi/2)^-,\ \theta_2\to(\pi/2)^+ \Rightarrow \phi(y)\to \pi.
	\end{align*}
	Hence
	\[
	f(0^+{+}iy)=\sqrt{1-y^2},\qquad f(0^-{+}iy)=-\sqrt{1-y^2}.
	\]
	\textbf{(ii) \(y>1\).} \(\theta_1=\theta_2=\pi/2\) on both sides, so \(\phi(y)=\pi/2\) and
	\[
	f(\pm0{+}iy)= i\sqrt{y^2-1}.
	\]
	\textbf{(iii) \(y<-1\).} \(\theta_1=\theta_2=-\pi/2\) (mod \(2\pi\)), so \(\phi(y)=-\pi/2\) and
	\[
	f(\pm0{+}iy)= -i\sqrt{y^2-1}.
	\]
	
	\subsection*{Branch cut and behaviour at infinity}
	The only jump occurs for \(|y|<1\) on the imaginary axis, hence the branch cut is
	\[
	\boxed{\,S=\{x+iy : x=0,\ |y|\le 1\}\, }.
	\]
	For \(|z|\to\infty\),
	\[
	f(z)=|z|e^{i\arg z}\,(1+o(1))=z\,(1+o(1)),
	\]
	so \(f(z)\sim z\). The image of \(S\) is the real segment \([-1,1]\); the imaginary axis outside \(S\) maps to the
	imaginary axis. Thus \(f:\mathbb C\setminus S\to \mathbb C\setminus[-1,1]\) is conformal with \(f(z)\sim z\).
	
	% ==================== Part (c) ====================
	\section*{4(c). A branch with two semi-infinite vertical cuts}
	Now define
	\[
	f(z)=(r_1r_2)^{1/2}e^{\,i(\theta_1+\theta_2)/2},\qquad
	-\frac{3\pi}{2}<\theta_1\le\frac{\pi}{2},\quad -\frac{\pi}{2}<\theta_2\le\frac{3\pi}{2}.
	\]
	
	\subsection*{Values across the imaginary axis}
	Let \(z=x+iy\), \(x\to0^\pm\).
	\begin{itemize}
		\item \(|y|<1\): we may keep \(\theta_1\approx -\pi/2\), \(\theta_2\approx \pi/2\) on both sides, hence
		\[
		f(\pm0{+}iy)=\sqrt{1-y^2}\quad(\text{no jump}).
		\]
		\item \(y>1\): crossing the axis forces \(\theta_1\) to jump by \(2\pi\); thus
		\[
		f(0^+{+}iy)=+i\sqrt{y^2-1},\qquad f(0^-{+}iy)=-i\sqrt{y^2-1}.
		\]
		\item \(y<-1\): analogously,
		\[
		f(0^+{+}iy)=-i\sqrt{y^2-1},\qquad f(0^-{+}iy)=+i\sqrt{y^2-1}.
		\]
	\end{itemize}
	Therefore the branch cut consists of the two rays \((-\infty i,-i]\cup[i,+\infty i)\).
	
	\subsection*{Values on the real axis and asymptotics by half-planes}
	For \(z=x\in\mathbb R\),
	\[
	r_1=r_2=\sqrt{x^2+1},\qquad
	\frac{\theta_1+\theta_2}{2}=\begin{cases}0,&x>0,\\ \pi,&x<0.\end{cases}
	\]
	Hence
	\[
	\boxed{\,f(x)=\operatorname{sgn}(x)\,\sqrt{x^2+1}\, }.
	\]
	As \(|z|\to\infty\),
	\[
	f(z)\sim |z|e^{i\arg z}\times
	\begin{cases}1,& \Re z>0,\\ -1,& \Re z<0,\end{cases}
	\]
	i.e.
	\[
	\boxed{\,f(z)\sim z\ \text{for }\Re z>0,\qquad f(z)\sim -z\ \text{for }\Re z<0.}
	\]
	
	
	\section*{Exact problem statement (as in the text)}
	
% ---------- Problem 4 (readable, line-broken version) ----------
\section*{4.\; Branches of \(\sqrt{z^{2}+1}\) — Problem statement}

\begin{enumerate}[label=(\alph*),leftmargin=2.1em]
	
	\item \textbf{(a)}\;
	Show that, if \(\zeta^{2}=z^{2}+1\) and
	\[
	z = i + r_{1}e^{i\theta_{1}} = -\,i + r_{2}e^{i\theta_{2}},
	\qquad r_{1},r_{2}>0,\ \theta_{1},\theta_{2}\in\mathbb{R},
	\]
	then
	\[
	\zeta \;=\; \pm\,(r_{1}r_{2})^{1/2}\,e^{\,i(\theta_{1}+\theta_{2})/2}.
	\]
	Explain briefly why \(z=\pm i\) are the branch points of the multifunction \((z^{2}+1)^{1/2}\).
	
	\medskip
	
	\item \textbf{(b)}\;
	Consider the branch defined by
	\[
	f(z) \;=\; (r_{1}r_{2})^{1/2} \, e^{\,i(\theta_{1}+\theta_{2})/2},
	\qquad -\tfrac{\pi}{2}<\theta_{1},\ \theta_{2}\le \tfrac{3\pi}{2}.
	\]
	State the values of \((\theta_{1}+\theta_{2})/2\) on either side of the imaginary axis,
	and hence compute \(f(\pm0+iy)\) in terms of \(y\) for \(y\in\mathbb{R}\).
	Deduce that the branch cut is
	\[
	S \;=\; \{\,x+iy:\ x=0,\ |y|\le 1\,\}.
	\]
	Show that \(f(z)\sim z\) as \(|z|\to\infty\)
	(i.e.\ \(f(z)/z\to 1\) as \(|z|\to\infty\)).
	Sketch the image of the cut \(z\)-plane \(\mathbb{C}\setminus S\) under the map \(\zeta=f(z)\).
	
	\medskip
	
	\item \textbf{(c)}\;
	Consider instead the branch
	\[
	f(z) \;=\; (r_{1}r_{2})^{1/2} \, e^{\,i(\theta_{1}+\theta_{2})/2},
	\qquad -\tfrac{3\pi}{2}<\theta_{1}\le \tfrac{\pi}{2},\quad
	-\tfrac{\pi}{2}<\theta_{2}\le \tfrac{3\pi}{2}.
	\]
	State the values of \((\theta_{1}+\theta_{2})/2\) on either side of the imaginary axis,
	and hence compute \(f(\pm0+iy)\) in terms of \(y\) for \(y\in\mathbb{R}\).
	Deduce that the branch cut is along the imaginary axis
	from \(z=-i\infty\) to \(z=-i\) and from \(z=i\) to \(z=i\infty\).
	Compute \(f(x)\) in terms of \(x\) for \(x\in\mathbb{R}\).
	Show that
	\[
	f(z)\sim
	\begin{cases}
		z,  & \text{as } |z|\to\infty \text{ with } \Re(z)>0,\\[2pt]
		-z,  & \text{as } |z|\to\infty \text{ with } \Re(z)<0.
	\end{cases}
	\]
	Sketch the images of the half-planes \(\Re(z)>0\) and \(\Re(z)<0\) under the map \(\zeta=f(z)\).
	
\end{enumerate}
	
	\bigskip\hrule\bigskip
	
	% =========================
	\section*{Solutions (with background)}
	
	Throughout, write
	\[
	z-i=r_{1}e^{i\theta_{1}},\qquad z+i=r_{2}e^{i\theta_{2}},\qquad r_{1},r_{2}>0,
	\]
	so that
	\[
	z^{2}+1=(z-i)(z+i)=r_{1}r_{2}\,e^{i(\theta_{1}+\theta_{2})}.
	\]
	
	\subsection*{(a) Formula for \(\zeta\) and branch points}
	Since \(\zeta^{2}=r_{1}r_{2}\,e^{i(\theta_{1}+\theta_{2})}\), any square root is
	\[
	\boxed{\ \zeta=\pm(r_{1}r_{2})^{1/2}e^{i(\theta_{1}+\theta_{2})/2}\ }.
	\]
	Branch points occur at zeros of \(z^{2}+1\), where an argument for a factor cannot be chosen continuously:
	these are exactly \(z=\pm i\).
	
	\subsection*{(b) Branch with finite vertical cut \(S=\{x=0,\,|y|\le 1\}\)}
	
	\textbf{Branch choice.} Take \(-\frac{\pi}{2}<\theta_{1},\theta_{2}\le \frac{3\pi}{2}\) and define
	\(f(z)=(r_{1}r_{2})^{1/2}e^{i(\theta_{1}+\theta_{2})/2}\).
	
	\textbf{Angles on the imaginary axis.} Let \(z=x+iy\) and denote
	\(\phi(y)=\tfrac12(\theta_{1}+\theta_{2})\).
	Use
	\[
	r_{1}=|y-1|,\quad r_{2}=|y+1|,\qquad
	\arg(i t)=\begin{cases}\pi/2,&t>0,\\ -\pi/2,&t<0.\end{cases}
	\]
	\emph{(i) \(|y|<1\).}
	As \(x\to 0^{+}\): \(\theta_{1}\to (-\pi/2)^{+}\), \(\theta_{2}\to (\pi/2)^{-}\Rightarrow \phi(y)\to 0\).
	As \(x\to 0^{-}\): to keep \(\theta_{1}\) in \((-\pi/2,3\pi/2]\) we add \(2\pi\),
	so \(\theta_{1}\to (3\pi/2)^{-}\), \(\theta_{2}\to(\pi/2)^{+}\Rightarrow \phi(y)\to \pi\).
	Hence
	\[
	\boxed{\,f(0^{+}+iy)=\sqrt{1-y^{2}},\qquad f(0^{-}+iy)=-\sqrt{1-y^{2}}\quad(|y|<1)\, }.
	\]
	
	\emph{(ii) \(y>1\).}
	\(\theta_{1}=\theta_{2}=\pi/2\) from either side, so
	\[
	\boxed{\,f(\pm 0+iy)= i\sqrt{y^{2}-1}\quad(y>1)\, }.
	\]
	
	\emph{(iii) \(y<-1\).}
	\(\theta_{1}=\theta_{2}=-\pi/2\) (mod \(2\pi\)), so
	\[
	\boxed{\,f(\pm 0+iy)= -i\sqrt{y^{2}-1}\quad(y<-1)\, }.
	\]
	
	\textbf{Branch cut.}
	The only sign jump occurs for \(|y|<1\) on the imaginary axis, hence
	\[
	\boxed{\,S=\{x+iy:\ x=0,\ |y|\le 1\}\, }.
	\]
	
	\textbf{Asymptotics.}
	As \(|z|\to\infty\),
	\[
	r_{1}r_{2}=|z-i||z+i|=|z^{2}+1|=|z|^{2}(1+o(1)),\qquad
	\theta_{1}+\theta_{2}=\arg(z^{2}+1)=2\arg z+o(1),
	\]
	so
	\[
	\boxed{\,f(z)\sim z\quad (|z|\to\infty)\, }.
	\]
	
	\textbf{Image sketch.}
	The cut segment \(S\) maps to the real interval \([-1,1]\) (opposite sides map to \((0,1]\) and \([-1,0)\));
	outside \(S\) the imaginary axis maps to the imaginary axis. Thus \(f:\mathbb C\setminus S\to\mathbb C\setminus[-1,1]\) is conformal with \(f(z)\sim z\).
	
	\subsection*{(c) Branch with two semi-infinite cuts}
	
	\textbf{Branch choice.}
	Take \(-\frac{3\pi}{2}<\theta_{1}\le \frac{\pi}{2}\) and \(-\frac{\pi}{2}<\theta_{2}\le \frac{3\pi}{2}\),
	and define \(f(z)=(r_{1}r_{2})^{1/2}e^{i(\theta_{1}+\theta_{2})/2}\).
	
	\textbf{Angles across the imaginary axis.}
	For \(z=x+iy\):
	
	\emph{(i) \(|y|<1\).}
	We may keep \(\theta_{1}\approx -\pi/2\) and \(\theta_{2}\approx \pi/2\) on both sides, hence
	\[
	\boxed{\,f(\pm 0+iy)=\sqrt{1-y^{2}}\quad(|y|<1)\, }.
	\]
	
	\emph{(ii) \(y>1\).}
	Crossing the axis forces \(\theta_{1}\) to jump by \(2\pi\); therefore
	\[
	\boxed{\,f(0^{+}+iy)=+i\sqrt{y^{2}-1},\qquad f(0^{-}+iy)=-i\sqrt{y^{2}-1}\quad(y>1)\, }.
	\]
	
	\emph{(iii) \(y<-1\).}
	Similarly,
	\[
	\boxed{\,f(0^{+}+iy)=-i\sqrt{y^{2}-1},\qquad f(0^{-}+iy)=+i\sqrt{y^{2}-1}\quad(y<-1)\, }.
	\]
	
	\textbf{Branch cuts.}
	The discontinuities occur on the rays
	\[
	\boxed{\,(-i\infty,-i]\ \cup\ [i,+i\infty)\, }.
	\]
	
	\textbf{Values on the real axis.}
	For \(z=x\in\mathbb R\), \(r_{1}=r_{2}=\sqrt{x^{2}+1}\) and
	\[
	\frac{\theta_{1}+\theta_{2}}{2}=
	\begin{cases}
		0,& x>0,\\
		\pi,& x<0.
	\end{cases}
	\qquad\Rightarrow\qquad
	\boxed{\,f(x)=\operatorname{sgn}(x)\,\sqrt{x^{2}+1}\, }.
	\]
	
	\textbf{Asymptotics by half-planes.}
	As \(|z|\to\infty\),
	\[
	\boxed{\,f(z)\sim z\ \ \text{for }\Re z>0,\qquad f(z)\sim -z\ \ \text{for }\Re z<0\, }.
	\]
	
	% =========================================================
	% Problem 5 — Images of sets under simple maps
	% (a) inversion, (b) sine of a vertical strip, (c) z+e^z on a horizontal strip
	% =========================================================
	\section{Problem 5: Images under $\zeta=\dfrac1z$, $\zeta=\sin z$, and $\zeta=z+e^{z}$}
	
	\subsection*{(a) Intersection of $|z-1|<1$ and $|z+i|<1$ under $\zeta=1/z$}
	
	\textbf{Background (inversion).}
	If $w=1/z$ and $|z-a|=|a|$ (a circle through $0$), then
	\[
	|1-aw|=|a||w|\ \Longleftrightarrow\ 1=2\Re(aw)\ \Longleftrightarrow\ \Re(aw)=\tfrac12,
	\]
	i.e. the image is a straight line.
	
	\textbf{Boundary images.}
	\[
	|z-1|=1 \ \longmapsto\ \Re(w)=\tfrac12,\qquad
	|z+i|=1 \ \longmapsto\ \Re((-i)w)=\tfrac12\ \Longleftrightarrow\ \Im(w)=\tfrac12 .
	\]
	
	\textbf{Which side is the interior?}
	\[
	|z-1|<1 \iff |1/w-1|<1 \iff |1-w|<|w|\iff \Re(w)>\tfrac12;
	\]
	\[
	|z+i|<1 \iff |1+iw|<|w|\iff \Im(w)>\tfrac12.
	\]
	
	\textbf{Image.} The lens $|z-1|<1\cap|z+i|<1$ maps to the quarter-plane
	\[
	\boxed{\ \{\,w:\ \Re w>\tfrac12,\ \Im w>\tfrac12\,\}\ }.
	\]
	
	\bigskip\hrule\bigskip
	
	\subsection*{(b) Strip $-\dfrac{\pi}{2}<x<\dfrac{\pi}{2}$ under $\zeta=\sin z$}
	
	Write $z=x+iy$. Then
	\[
	\sin z=\sin x\,\cosh y \;+\; i\,\cos x\,\sinh y.
	\]
	
	\textbf{Boundary lines.}
	\[
	\sin\!\Big(\pm \frac{\pi}{2}+iy\Big)=\pm \cosh y\in\mathbb{R},\quad |\cosh y|\ge1.
	\]
	Hence the boundaries map onto the real rays $[1,\infty)$ and $(-\infty,-1]$.
	
	\textbf{Conformality.}
	$\cos z\neq 0$ on $|x|<\pi/2$ (zeros lie on $x=\pi/2+k\pi$), so $\sin$ is conformal and injective on the strip.
	
	\textbf{Image.} The image is the simply connected domain whose boundary is those rays; by symmetry (the line $x=0$ maps to the entire imaginary axis),
	\[
	\boxed{\ \sin:\ \{-\tfrac{\pi}{2}<\Re z<\tfrac{\pi}{2}\}\ \xrightarrow[\text{1--1}]{}
		\ \mathbb{C}\setminus\big((-\infty,-1]\cup[1,\infty)\big). \ }
	\]
	
	\bigskip\hrule\bigskip
	
	\subsection*{(c) Strip $-\pi<y<\pi$ under $\zeta=f(z)=z+e^{z}$}
	
	Let $z=x+iy$. Then
	\[
	f(z)=\big(x+e^{x}\cos y\big)\;+\;i\big(y+e^{x}\sin y\big).
	\]
	
	\textbf{Critical points (hint).}
	\[
	f'(z)=1+e^{z}=0 \iff e^{z}=-1 \iff z=i(2k+1)\pi .
	\]
	There are no critical points in the open strip $-\pi<y<\pi$; the boundary points $z=\pm i\pi$ are the only adjacent critical points.
	
	\textbf{Boundary images.}
	For $y=\pm\pi$, $\sin y=0$, $\cos y=-1$:
	\[
	f(x\pm i\pi)=\big(x-e^{x}\big)\ \pm\ i\pi.
	\]
	Here $u(x)=x-e^{x}$ has a unique maximum $-1$ at $x=0$ and $u(x)\to-\infty$ as $x\to\pm\infty$.
	Thus each boundary maps to a horizontal half-line
	\[
	(-\infty,-1]\,+\,i\pi,\qquad (-\infty,-1]\,-\,i\pi,
	\]
	meeting at the critical values $-1\pm i\pi$.
	
	\textbf{Interior and image.}
	With no interior critical points, $f$ is conformal on the strip.
	For fixed $y\in(-\pi,\pi)$, $f(x+iy)\to -\infty+iy$ as $x\to-\infty$ and
	$f(x+iy)\sim e^{x}e^{iy}\to\infty$ along the ray $\arg\zeta=y$ as $x\to+\infty$.
	Hence the image is the plane with the two horizontal half-lines removed:
	\[
	\boxed{\ \mathbb{C}\setminus\Big((-\infty,-1]+i\pi\ \cup\ (-\infty,-1]-i\pi\Big). \ }
	\]
	
	\section{Problem 6 — Conformal mappings}
	
	\subsection*{Exact problem statement}
	
	\noindent
	(a)\; $D$ is the region exterior to the two circles $|z-1|=1$ and $|z+1|=1$.
	Find a conformal mapping from $D$ to the exterior of the unit circle.\\
	\emph{Hint:} First apply an inversion with respect to the origin $(z\mapsto 1/z)$,
	then rotate and scale so as to use the exponential function to get to a half plane,
	from where use Möbius.
	
	\medskip
	\noindent
	(b)\; Find a conformal map from the unit disc $|z|<1$ onto the strip $-\pi/2<\eta<\pi/2$,
	taking the origin to the origin, $1$ to $\xi=+\infty$, $-1$ to $\xi=-\infty$.
	(Thus, the upper half of the unit circle maps to $\eta=\pi/2$ and the lower half to $\eta=-\pi/2$.)
	
	\medskip
	\noindent
	(c)\; Find a conformal map of the quarter-disc $0<|z|<1,\ 0<\arg z<\pi/2$
	to the upper half-plane $\eta>0$, taking $0$ to $0$, $1$ to $1$, and $i$ to $\infty$.
	
	\bigskip\hrule\bigskip
	
	\subsection*{Solutions (step by step)}
	
	\paragraph{(a) Exterior of the two circles $\to$ exterior of $|\zeta|=1$.}
	\begin{itemize}
		\setlength{\itemsep}{0.4em}
		\item \textbf{Invert.} Let $w=1/z$. Then the circles through $0$ become lines:
		\[
		|z-1|=1 \Rightarrow \Re w=\tfrac12,\qquad
		|z+1|=1 \Rightarrow \Re w=-\tfrac12.
		\]
		Hence $D$ maps to the vertical strip $S=\{-\tfrac12<\Re w<\tfrac12\}$.
		\item \textbf{Rotate/scale to a horizontal strip.} Put $s=i\pi w$.
		Then $|\Im s|<\pi/2$.
		\item \textbf{Exponential to a half-plane.} Let $u=e^{\,s}$.
		Then $|\Im s|<\pi/2 \;\Longrightarrow\; \Re u>0$.
		\item \textbf{Möbius to exterior unit disk.} Use
	% needs \usepackage{amsmath}
	\[
	\Phi(u)=\frac{u+1}{u-1},\qquad \operatorname{Re} u>0 \longmapsto |\zeta|>1.
	\]
	
	\end{itemize}
	Therefore a concrete map is
	\[
	\boxed{\ \zeta=\Phi\!\big(e^{\,i\pi/z}\big)=\dfrac{e^{\,i\pi/z}+1}{e^{\,i\pi/z}-1}\ }.
	\]
	
	\bigskip
	
	\paragraph{(b) Unit disk $\to$ strip $-\pi/2<\Im\zeta<\pi/2$.}
	\begin{itemize}
		\setlength{\itemsep}{0.4em}
		\item \textbf{Disk to right half-plane (Cayley).}
		\[
		w=\frac{1+z}{1-z},\qquad |z|<1 \mapsto \Re w>0,
		\]
		and $z=0\mapsto w=1$, $z=1\mapsto w=\infty$, $z=-1\mapsto w=0$.
		\item \textbf{Log to a horizontal strip.}
		\[
		\zeta=\log w \quad (\text{principal branch}),\qquad \Im\zeta\in\big(-\tfrac{\pi}{2},\tfrac{\pi}{2}\big).
		\]
		Then $\zeta(0)=0$, $\zeta(1)=+\infty$, $\zeta(-1)=-\infty$.
	\end{itemize}
	Hence
	\[
	\boxed{\ \zeta=\log\!\left(\frac{1+z}{1-z}\right)\ }.
	\]
	
	\bigskip
	
	\paragraph{(c) Quarter-disk $\to$ upper half-plane; $0\mapsto0$, $1\mapsto1$, $i\mapsto\infty$.}
	\begin{itemize}
		\setlength{\itemsep}{0.4em}
		\item \textbf{Square to upper semicircle.}
		\[
		t=z^{2}:\quad \{0<|z|<1,\ 0<\arg z<\tfrac{\pi}{2}\}
		\longmapsto \{|t|<1,\ \Im t>0\}.
		\]
		Moreover $0\mapsto0$, $1\mapsto1$, $i\mapsto-1$.
		\item \textbf{Real Möbius to upper half-plane.}
		Find $M(t)$ with $M(-1)=\infty$, $M(0)=0$, $M(1)=1$.
		A suitable choice is
		\[
		M(t)=\frac{2t}{t+1}.
		\]
	\end{itemize}
	Therefore
	\[
	\boxed{\ \zeta=M(z^{2})=\frac{2z^{2}}{z^{2}+1}\ }.
	\]
	It maps the quarter-disk conformally onto $\Im\zeta>0$ and has
	$\zeta(0)=0$, $\zeta(1)=1$, $\zeta(i)=\infty$.
	
	
	% =====================================================================
	% Problem 6(c) — Quarter-disk to upper half-plane: full explanation
	% =====================================================================
	\section*{6(c)\quad Quarter-disk $0<|z|<1,\ 0<\arg z<\tfrac{\pi}{2}$
		$\longrightarrow$ upper half-plane $\Im\zeta>0$,
		with $0\mapsto0$, $1\mapsto1$, $i\mapsto\infty$}
	
	\subsection*{Goal}
	Construct an explicit conformal map $\zeta=\zeta(z)$ that maps the open quarter-disk
	\[
	Q=\{\,z:\ 0<|z|<1,\ 0<\arg z<\tfrac{\pi}{2}\,\}
	\]
	onto the upper half-plane $\mathbb H=\{\zeta:\ \Im\zeta>0\}$, sending the three marked boundary
	points
	\[
	z=0 \longmapsto \zeta=0,\qquad
	z=1 \longmapsto \zeta=1,\qquad
	z=i \longmapsto \zeta=\infty .
	\]
	
	\subsection*{Step 1: Square to unfold the angle}
	Let
	\[
	t=z^{2}.
	\]
	Write $z=re^{i\theta}$ with $0<r<1$ and $0<\theta<\frac{\pi}{2}$. Then
	\[
	t=r^{2}e^{i(2\theta)}\quad\Rightarrow\quad |t|<1,\ \ 0<\arg t<\pi \iff \Im t>0.
	\]
	Therefore the squaring map $z\mapsto t=z^{2}$ sends $Q$ \emph{bijectively and conformally}
	onto the \emph{open upper semicircle}
	\[
	S=\{\,t:\ |t|<1,\ \Im t>0\,\}.
	\]
	On the boundary rays/arcs we have
	\[
	0\mapsto 0,\qquad 1\mapsto 1,\qquad i\mapsto -1.
	\]
	(Conformality is immediate, since $(z^{2})'=2z\neq 0$ on $Q$.)
	
	\subsection*{Step 2: Real Möbius map fixing the three marked points and sending $S$ to $\mathbb H$}
	Seek a real-coefficient Möbius transformation
	\[
	M(t)=\frac{a t+b}{c t+d},\qquad a,b,c,d\in\mathbb R,\ \ ad-bc\neq0,
	\]
	such that
	\[
	M(-1)=\infty,\qquad M(0)=0,\qquad M(1)=1.
	\]
	Imposing the conditions:
	
	\begin{itemize}
		\item $M(-1)=\infty \ \Rightarrow\ c(-1)+d=0 \ \Rightarrow\ d=c$.
		\item $M(0)=0 \ \Rightarrow\ b/d=0 \ \Rightarrow\ b=0$.
		\item $M(1)=1 \ \Rightarrow\ a/(c+d)=1 \ \Rightarrow\ a=c+d=2c$.
	\end{itemize}
	
	Up to an overall nonzero constant (which does not change the map), we can take $c=1$ and obtain
	\[
	\boxed{\,M(t)=\frac{2t}{t+1}\,}.
	\]
	
	\paragraph{Why $M$ maps $S$ into $\Im\zeta>0$ (and in fact maps the whole upper half-plane to itself).}
	Let $t=x+iy$ with $y>0$. Then
	\[
	M(t)=\frac{2t}{1+t},\qquad
	\Im M(t)=\frac{2\,\Im\big(t(1+\overline{t})\big)}{|1+t|^{2}}
	=\frac{2\,\Im t}{|1+t|^{2}}
	=\frac{2y}{|1+t|^{2}}
	>0.
	\]
	Hence $M$ maps $\{\,\Im t>0\,\}$ into itself; in particular it maps the upper semicircle $S$
	(which lies in $\Im t>0$) into the upper half-plane. Moreover
	\[
	M'(t)=\frac{2}{(1+t)^{2}}\neq0\quad\text{on}\ S,
	\]
	so $M$ is conformal (locally injective) there. The three-point normalization guarantees that it
	sends the three distinguished boundary points as prescribed:
	\[
	M(-1)=\infty,\qquad M(0)=0,\qquad M(1)=1.
	\]
	
	\subsection*{Step 3: Compose and verify}
	Define the desired map
	\[
	\boxed{\ \zeta(z)=M\big(z^{2}\big)=\frac{2z^{2}}{z^{2}+1}\ }.
	\]
	Then the normalization follows immediately:
	\[
	\zeta(0)=0,\qquad \zeta(1)=1,\qquad \zeta(i)=\infty .
	\]
	Because $t=z^{2}$ has $\Im t>0$ on $Q$ and $\Im M(t)>0$ for $\Im t>0$, we have
	\[
	\Im \zeta(z)=\Im\big(M(z^{2})\big)>0\quad\text{for all }z\in Q.
	\]
	Also $\zeta'(z)=M'(z^{2})\cdot 2z=\dfrac{4z}{(1+z^{2})^{2}}\neq0$ on $Q$, hence the map is conformal
	throughout the domain.
	
	\subsection*{Boundary correspondence (for intuition)}
	\begin{itemize}
		\item The \emph{real radius} segment $\{\,0<z<1\,\}$ maps by $t=z^{2}$ to $(0,1)$ and then by $M$
		to $(0,1)\subset\mathbb R$; thus that edge lands on the real axis between $0$ and $1$.
		\item The \emph{imaginary radius} segment $\{\,0<iz<i\,\}$ maps by $t=z^{2}$ to $(0,-1)$ and then by $M$
		to $(0,\infty)$ on $\mathbb R$; thus the edge $0\to i$ lands on the ray $[0,\infty)$.
		\item The \emph{circular arc} $\{\,e^{i\theta}: 0<\theta<\tfrac{\pi}{2}\,\}$ maps by $t=z^{2}$ to the unit
		semicircle $\{\,e^{i\phi}: 0<\phi<\pi\,\}$ and then by $M$ to the real line as a boundary set.
	\end{itemize}
	The interior of the quarter-disk maps onto $\Im\zeta>0$ by the positivity computation for $\Im M(t)$.
	
	\subsection*{Conclusion}
	The composition
	\[
	\boxed{\ \zeta(z)=\frac{2z^{2}}{z^{2}+1}\ }
	\]
	is a conformal bijection from the open quarter-disk $Q$ onto the upper half-plane $\mathbb H$,
	satisfying the three mapping conditions $0\mapsto0$, $1\mapsto1$, $i\mapsto\infty$.
	
	
	% ===== Problem 6(b): unit disk to strip — how & why =====
	
	\subsection*{(b) Unit disk $\to$ strip $-\dfrac{\pi}{2}<\Im\zeta<\dfrac{\pi}{2}$}
	
	\textbf{Disk to right half-plane (Cayley).}
	\[
	w=\frac{1+z}{1-z},\qquad |z|<1 \;\Longrightarrow\; \operatorname{Re} w>0,
	\]
	and on the marked points
	\[
	z=0 \mapsto w=1,\qquad z=1 \mapsto w=\infty,\qquad z=-1 \mapsto w=0.
	\]
	Moreover, the unit circle $|z|=1\setminus\{1\}$ maps to the imaginary axis $\operatorname{Re} w=0$.
	
	\medskip
	\textbf{Log to a horizontal strip.}
	Use the principal logarithm on the right half-plane:
	\[
	\zeta=\log w=\ln|w|+i\,\arg(w),\qquad \arg(w)\in\big(-\tfrac{\pi}{2},\tfrac{\pi}{2}\big),
	\]
	so
	\[
	-\frac{\pi}{2} < \Im \zeta < \frac{\pi}{2}.
	\]
	At the marked points:
	\[
	\zeta(0)=\log 1=0,\qquad
	\zeta(1)=\log(\infty)=+\infty,\qquad
	\zeta(-1)=\log(0)=-\infty.
	\]
	
	\medskip
	\textbf{Which boundary goes where.}
	For $z=e^{i\theta}$,
	\[
	w=\frac{1+e^{i\theta}}{1-e^{i\theta}}
	= -\,i\,\cot\!\left(\frac{\theta}{2}\right)\in i\mathbb{R}.
	\]
	Hence
	\[
	0<\theta<\pi\ (\text{upper semicircle}) \Rightarrow \Im w<0 \Rightarrow \Im \zeta=-\frac{\pi}{2},
	\]
	\[
	\pi<\theta<2\pi\ (\text{lower semicircle}) \Rightarrow \Im w>0 \Rightarrow \Im \zeta=+\frac{\pi}{2}.
	\]
	(If one wants the upper semicircle to land on $\Im\zeta=+\pi/2$, use $\zeta=-\log\!\left(\frac{1+z}{1-z}\right)$.)
	
	\medskip
	\textbf{Hence}
	\[
	\boxed{\ \zeta=\log\!\left(\frac{1+z}{1-z}\right)\ }
	\]
	is a conformal bijection from $|z|<1$ onto the strip $-\pi/2<\Im\zeta<\pi/2$ with
	$0\mapsto 0$, $1\mapsto +\infty$, $-1\mapsto -\infty$.
	
	
	% =====================================================================
	% Problem 6(a) — Exterior of two tangent circles -> Exterior of unit disk
	% Full, non-shortened LaTeX with detailed steps and checks
	% =====================================================================
	
	\subsection*{(a) Exterior of the two circles $\boldsymbol{|z-1|=1}$ and $\boldsymbol{|z+1|=1}$ 
		$\;\longrightarrow\;$ exterior of $\boldsymbol{|\zeta|=1}$}
	
	\paragraph{Goal.}
	Let
	\[
	D=\big\{z\in\mathbb C:\ |z-1|\ge 1,\ |z+1|\ge 1\big\}\setminus\big(\{|z-1|=1\}\cup\{|z+1|=1\}\big),
	\]
	the region \emph{exterior to both} unit circles centered at $\pm1$.
	Construct an explicit conformal bijection $F:D\to\{\zeta:\ |\zeta|>1\}$.
	
	\paragraph{Step 1: Inversion.}
	Let $w=\dfrac{1}{z}$.
	For any $a\in\mathbb C$, the circle $\{|z-a|=|a|\}$ (which passes through $0$) satisfies
	\[
	|z-a|=|a|
	\iff \left|\frac{1}{w}-a\right|=|a|
	\iff |1-aw|=|a||w|
	\iff 1=2\,\Re(aw).
	\]
	Hence circles through $0$ invert to \emph{lines}.
	For $a=1$ and $a=-1$ we obtain
	\[
	|z-1|=1 \iff \Re w=\frac{1}{2},
	\qquad
	|z+1|=1 \iff \Re w=-\frac{1}{2}.
	\]
	
	\emph{Which side is the exterior?}  Using $|1-aw|^{2}=(1-2\Re(aw)+|a|^{2}|w|^{2})$ we compute
	\[
	|z-1|>1
	\iff |1-w|>|w|
	\iff 1-2\Re w>0
	\iff \Re w<\frac{1}{2},
	\]
	\[
	|z+1|>1
	\iff |1+ w|>|w|
	\iff 1+2\Re w>0
	\iff \Re w>-\frac{1}{2}.
	\]
	Therefore the inversion $w=1/z$ sends $D$ bijectively onto the \emph{vertical strip}
	\[
	S=\Big\{\,w\in\mathbb C:\ -\tfrac12<\Re w<\tfrac12\,\Big\}.
	\]
	(Since $z\neq 0$ on $D$, the inversion is analytic and injective there.)
	
	\paragraph{Step 2: Rotate/scale to a horizontal strip.}
	Define
	\[
	s=i\pi w.
	\]
	Write $w=u+iv$ with $-\tfrac12<u<\tfrac12$. Then $s=i\pi(u+iv)=-\pi v+i\pi u$, so
	\[
	|\Im s|=|\pi u|<\frac{\pi}{2}.
	\]
	Consequently $S$ is mapped conformally onto the horizontal strip
	\[
	T=\Big\{\,s\in\mathbb C:\ |\Im s|<\frac{\pi}{2}\,\Big\}.
	\]
	
	\paragraph{Step 3: Exponential to the right half-plane.}
	Set
	\[
	u=e^{\,s}.
	\]
	If $s=\sigma+i\tau$ with $|\tau|<\pi/2$, then
	\[
	\Re u = \Re\big(e^{\sigma}(\cos\tau+i\sin\tau)\big)=e^{\sigma}\cos\tau>0,
	\]
	because $\cos\tau>0$ on $(-\pi/2,\pi/2)$.
	Hence $e^{(\cdot)}$ maps $T$ conformally and bijectively onto the \emph{right half-plane}
	\[
	H=\{\,u\in\mathbb C:\ \Re u>0\,\},
	\]
	and the boundary lines $\Im s=\pm\pi/2$ go to the imaginary axis $\Re u=0$.
	
	\paragraph{Step 4: Möbius map from the right half-plane to the exterior disk.}
	Consider
	\[
	\Phi(u)=\frac{u+1}{u-1}.
	\]
	For $u\in H$ we have the equivalence
	\[
	|\Phi(u)|>1
	\iff |u+1|>|u-1|
	\iff (u+1)(\bar u+1)>(u-1)(\bar u-1)
	\iff 4\,\Re u>0,
	\]
	which holds precisely on $H$. Thus $\Phi$ maps $H$ conformally onto $\{\zeta:\ |\zeta|>1\}$ and
	sends the imaginary axis to the unit circle $|\zeta|=1$.
	
	\paragraph{Step 5: Composition and explicit formula.}
	The composition of the four conformal bijections yields
	\[
	F:D \xrightarrow{\,z\mapsto w=1/z\,} S
	\xrightarrow{\,w\mapsto s=i\pi w\,} T
	\xrightarrow{\,s\mapsto u=e^{\,s}\,} H
	\xrightarrow{\,u\mapsto \Phi(u)\,} \{|\zeta|>1\}.
	\]
	Therefore an explicit mapping is
	\[
	\boxed{\quad \zeta = F(z) = \Phi\!\big(e^{\,i\pi/z}\big)
		= \frac{e^{\,i\pi/z}+1}{\,e^{\,i\pi/z}-1\,}\quad }.
	\]
	Every factor is analytic and injective on the relevant domain; hence $F$ is conformal and bijective from $D$ onto $\{|\zeta|>1\}$.
	
	\paragraph{Asymptotic check at infinity (sanity).}
	As $z\to\infty$ we have $e^{\,i\pi/z}=1+\dfrac{i\pi}{z}+O(z^{-2})$, so
	\[
	\zeta
	=\frac{2+\dfrac{i\pi}{z}+O(z^{-2})}{\dfrac{i\pi}{z}+O(z^{-2})}
	= \frac{2}{i\pi}\,z\,\Big(1+O(z^{-1})\Big)
	= -\,\frac{2i}{\pi}\,z\Big(1+O(z^{-1})\Big),
	\]
	which confirms that far out in $D$ the map behaves like an affine scaling of $z$, as expected.
	
	\paragraph{Boundary correspondence (orientation).}
	\begin{itemize}
		\item The line $\Re w=\tfrac12$ (image of $|z-1|=1$) becomes $\Im s=\tfrac{\pi}{2}$,
		then the imaginary axis, and finally the unit circle $|\zeta|=1$.
		\item The line $\Re w=-\tfrac12$ (image of $|z+1|=1$) becomes $\Im s=-\tfrac{\pi}{2}$,
		then the same imaginary axis, and again $|\zeta|=1$.
		\item The interior $-\tfrac12<\Re w<\tfrac12$ maps to $|\Im s|<\tfrac{\pi}{2}$,
		then to $\Re u>0$, and finally to $|\zeta|>1$.
	\end{itemize}
	This completes the construction and detailed verification.
	
	% =====================================================================
	% Problem 7 — Half–plane with a circular hole: map to an annulus and solve
	% DO NOT SHORTEN — full statement, background, derivation, and checks
	% =====================================================================
	
	\section{Problem 7: Temperature in a right half–plane with a removed disk}
	
	\subsection*{Exact statement (verbatim)}
	The domain $D$ consists of the right-hand half plane $x>0$ with the circle $|z-a|=b$, $0<b<a$, and its interior removed.
	Find the temperature $u(x,y)$ in steady heat flow if $u=0$ on the $y$ axis, $u=1$ on $|z-a|=b$, and $u\to0$ at infinity.
	
	\medskip
	\noindent\textit{Hint:} Show that the mapping $\displaystyle \zeta=\frac{z-\alpha}{z+\alpha}$, with $\alpha$ real and positive,
	takes $D$ onto an annular region with the imaginary axis mapping to $|\zeta|=1$ and show that, if $\alpha^{2}=a^{2}-b^{2}$, then the image of $D$ is a concentric circular annulus.
	
	\bigskip\hrule\bigskip
	
	\subsection*{Background: why the suggested Möbius map produces an annulus}
	Consider the real-parameter Möbius map
	\[
	\zeta=\zeta(z)=\frac{z-\alpha}{z+\alpha}, \qquad \alpha>0.
	\]
	\paragraph{Imaginary axis.}
	For $z=iy$ we have
	\[
	\left|\frac{iy-\alpha}{iy+\alpha}\right|=1 \quad\text{(numerator and denominator are conjugates),}
	\]
	so the $y$–axis maps to the unit circle $|\zeta|=1$.
	
	\paragraph{Image of the circle $\boldsymbol{|z-a|=b}$.}
	Write $z=a+be^{i\theta}$ and compute
	\[
	|\zeta|^{2}
	=\frac{|a-\alpha+be^{i\theta}|^{2}}{|a+\alpha+be^{i\theta}|^{2}}
	=\frac{A+B\cos\theta}{C+D\cos\theta},
	\]
	where
	\[
	A=(a-\alpha)^{2}+b^{2}, \qquad B=2b(a-\alpha), \qquad
	C=(a+\alpha)^{2}+b^{2}, \qquad D=2b(a+\alpha).
	\]
	For $|\zeta|$ to be constant on the circle (i.e. the image is a circle \emph{centered at the origin}),
	the $\cos\theta$–dependence must drop out:
	\[
	\frac{A+B\cos\theta}{\,C+D\cos\theta\,}\equiv\frac{A}{C}
	\quad\Longleftrightarrow\quad AD=BC.
	\]
	We now verify that
	\[
	AD-BC
	=2b\Big[(a+\alpha)((a-\alpha)^{2}+b^{2})-(a-\alpha)((a+\alpha)^{2}+b^{2})\Big]
	=4b\alpha\,(\alpha^{2}+b^{2}-a^{2}).
	\]
	Thus $AD=BC$ is equivalent to
	\[
	\boxed{\ \alpha^{2}=a^{2}-b^{2}\ }.
	\]
	With this choice (note $a>b>0\Rightarrow \alpha>0$), the ratio becomes \emph{constant}:
	\[
	|\zeta|^{2}=\frac{A}{C}
	=\frac{(a-\alpha)^{2}+b^{2}}{(a+\alpha)^{2}+b^{2}}
	=\frac{2a(a-\alpha)}{2a(a+\alpha)}
	=\frac{a-\alpha}{a+\alpha}.
	\]
	Hence the circle $|z-a|=b$ maps to the circle $|\zeta|=\rho$ with
	\[
	\boxed{\ \rho=\sqrt{\frac{a-\alpha}{a+\alpha}}\in(0,1). \ }
	\]
	
	\paragraph{Image of the domain.}
	Therefore, for $\alpha=\sqrt{a^{2}-b^{2}}$, the map $z\mapsto\zeta$ sends the domain $D$
	(right half–plane with the closed disk $\{|z-a|\le b\}$ removed) onto the concentric circular \emph{annulus}
	\[
	\mathcal A=\{\ \rho<|\zeta|<1\ \}.
	\]
	The boundary pieces correspond as follows:
	\[
	\{x=0\}\ \longmapsto\ |\zeta|=1,
	\qquad
	\{|z-a|=b\}\ \longmapsto\ |\zeta|=\rho.
	\]
	
	\bigskip\hrule\bigskip
	
	\subsection*{Solving the Dirichlet problem on the annulus}
	Let $U(\zeta)$ denote the temperature in the $\zeta$–plane. Harmonicity is preserved by conformal maps, so
	$U$ is harmonic on $\mathcal A$ and the boundary data become
	\[
	U=0 \ \ \text{on}\ \ |\zeta|=1,
	\qquad
	U=1 \ \ \text{on}\ \ |\zeta|=\rho.
	\]
	By rotational symmetry, $U$ depends only on $r=|\zeta|$. The general radial harmonic function on an annulus is
	$A\log r + B$. Imposing the boundary values gives
	\[
	U(r)=\frac{\log(1/r)}{\log(1/\rho)}
	=\frac{-\log r}{-\log \rho}
	=\frac{\log r}{\log \rho}
	\quad\text{for}\quad \rho<r<1.
	\]
	(Any of the three displayed forms is correct; we will use the first one momentarily.)
	
	\bigskip\hrule\bigskip
	
	\subsection*{Pullback to the original $z$–plane and final formula}
	The physical temperature is $u(z)=U(\zeta(z))$ with $r=|\zeta(z)|=\left|\frac{z-\alpha}{z+\alpha}\right|$.
	Hence
	\[
	u(z)
	=U\!\left(\left|\frac{z-\alpha}{z+\alpha}\right|\right)
	=\frac{\displaystyle \log\!\left(\frac{1}{\left|\frac{z-\alpha}{z+\alpha}\right|}\right)}
	{\displaystyle \log\!\left(\frac{1}{\rho}\right)}
	=\frac{\displaystyle \log\left|\frac{z+\alpha}{z-\alpha}\right|}
	{\displaystyle \log\!\left(\frac{1}{\rho}\right)}.
	\]
	Since $\displaystyle \rho=\sqrt{\frac{a-\alpha}{a+\alpha}}$ we have
	\[
	\log\!\left(\frac{1}{\rho}\right)
	=\frac{1}{2}\,\log\!\left(\frac{a+\alpha}{a-\alpha}\right).
	\]
	Therefore a convenient explicit form is
	\[
	\boxed{\quad
		u(z)
		=\frac{2\,\log\!\left|\dfrac{z+\alpha}{z-\alpha}\right|}
		{\log\!\left(\dfrac{a+\alpha}{a-\alpha}\right)},
		\qquad
		\alpha=\sqrt{a^{2}-b^{2}}.
		\quad}
	\]
	
	\bigskip
	
	\subsection*{Checks of the boundary conditions and far field}
	\begin{itemize}
		\item \textbf{On the $y$--axis $x=0$:} for $z=iy$ we have $\left|\dfrac{z-\alpha}{z+\alpha}\right|=1$,
		so the numerator $\log\left|\dfrac{z+\alpha}{z-\alpha}\right|=0$ and hence $u=0$.
		\item \textbf{On the circle $|z-a|=b$:} by construction $|\zeta|=\rho$, i.e.
		$\left|\dfrac{z-\alpha}{z+\alpha}\right|=\rho$, so
		$\log\left|\dfrac{z+\alpha}{z-\alpha}\right|=\log\!\left(\dfrac{1}{\rho}\right)$
		and the fraction evaluates to $u=1$.
		\item \textbf{At infinity:} as $|z|\to\infty$, $\left|\dfrac{z+\alpha}{z-\alpha}\right|\to 1$,
		so the numerator $\to 0$ and $u(z)\to 0$.
	\end{itemize}
	
	\bigskip\hrule\bigskip
	
	\subsection*{Summary}
	With $\alpha=\sqrt{a^{2}-b^{2}}$ and $\displaystyle \rho=\sqrt{\frac{a-\alpha}{a+\alpha}}$,
	the Möbius map $\displaystyle \zeta=\frac{z-\alpha}{z+\alpha}$ carries
	$D$ onto the annulus $\{\rho<|\zeta|<1\}$.
	The unique harmonic function with $U=1$ on $|\zeta|=\rho$ and $U=0$ on $|\zeta|=1$
	is $U(r)=\log(1/r)/\log(1/\rho)$, and hence
	\[
	u(z)=\frac{2\,\log\!\left|\dfrac{z+\alpha}{z-\alpha}\right|}
	{\log\!\left(\dfrac{a+\alpha}{a-\alpha}\right)}.
	\]
	This $u$ satisfies all the prescribed boundary conditions and $u(z)\to 0$ as $|z|\to\infty$.
	

\section*{Hyperbolic maps: $\sinh$, $\cosh$, $\tanh$, $\coth$ — background, zeros/poles, and strip mappings}

\subsection*{1.\ $w=\sinh z$}
\textbf{Background:} entire, odd; period $2\pi i$. \quad
\textbf{Zeros:} $z=i\pi k$; \textbf{poles:} none.\\
\textbf{Decomposition: } $w=u+iv=\sinh x\cos y + i\,\cosh x\sin y$.\\
\textbf{Vertical lines } $x=\text{const}$: $\left(\frac{u}{\sinh x}\right)^2+\left(\frac{v}{\cosh x}\right)^2=1$ (ellipses, foci $\pm i$).\\
\textbf{Horizontal lines } $y=\text{const}$: $\left(\frac{u}{\cos y}\right)^2-\left(\frac{v}{\sin y}\right)^2=-1$ (hyperbolas, foci $\pm i$).\\
\textbf{Strip mapping:}
\[
\boxed{\,\sinh:\ \{|\Im z|<\tfrac{\pi}{2}\}\ \xrightarrow{\ 1\text{--}1\ }\ \mathbb C\setminus i\big((-\infty,-1]\cup[1,\infty)\big)\,}
\]
Boundaries $y=\pm\frac{\pi}{2}$ map to imaginary rays $\pm i[1,\infty)$.

\bigskip\hrule\bigskip

\subsection*{2.\ $w=\cosh z$}
\textbf{Background:} entire, even; period $2\pi i$. \quad
\textbf{Zeros:} $z=i\pi(k+\frac12)$; \textbf{poles:} none.\\
\textbf{Decomposition: } $w=u+iv=\cosh x\cos y + i\,\sinh x\sin y$.\\
\textbf{Vertical lines } $x=\text{const}$: $\left(\frac{u}{\cosh x}\right)^2+\left(\frac{v}{\sinh x}\right)^2=1$ (ellipses, foci $\pm1$).\\
\textbf{Horizontal lines } $y=\text{const}$: $\left(\frac{u}{\cos y}\right)^2-\left(\frac{v}{\sin y}\right)^2=1$ (hyperbolas, foci $\pm1$).\\
\textbf{Strip mapping:}
\[
\boxed{\,\cosh:\ \{|\Im z|<\pi\}\ \xrightarrow{\ 2\text{--}1\ }\ \mathbb C\setminus\big((-\infty,-1]\cup[1,\infty)\big)\,}
\]
On the half-strip $0<|\Im z|<\pi$, the map is 1--1 onto the slit plane.

\bigskip\hrule\bigskip

\subsection*{3.\ $w=\tanh z$}
\textbf{Background:} meromorphic, odd; period $i\pi$. \quad
\textbf{Zeros:} $z=i\pi k$; \textbf{poles:} $z=i\pi(k+\tfrac12)$.\\
\textbf{Strip $\to$ disk:}
\[
\boxed{\,\tanh:\ \{|\Im z|<\tfrac{\pi}{2}\}\ \xrightarrow{\ 1\text{--}1\ }\ \mathbb D\,}
\]
Real line maps to $(-1,1)$; as $|\Im z|\to\pi/2$, $|\tanh z|\to1$.

\bigskip\hrule\bigskip

\subsection*{4.\ $w=\coth z$}
\textbf{Background:} meromorphic, odd; period $i\pi$. \quad
\textbf{Zeros:} $z=i\pi(k+\tfrac12)$; \textbf{poles:} $z=i\pi k$.\\
\textbf{Strip $\to$ exterior disk:}
\[
\boxed{\,\coth:\ \{|\Im z|<\tfrac{\pi}{2}\}\setminus\{0\}\ \xrightarrow{\ 1\text{--}1\ }\ \{\,|w|>1\,\}\,}
\]
Real line maps to $(-\infty,-1)\cup(1,\infty)$; $|\Im z|\to\pi/2$ gives $|w|\to1^+$.



% =====================================================================
% Hyperbolic Maps: sinh, cosh, tanh, coth — theory, geometry, strip mappings
% Full LaTeX: background, zeros/poles, geometry of lines, strip/band maps,
% inverse branches, boundary correspondence, and useful identities
% =====================================================================

\section*{Hyperbolic Conformal Maps: $\sinh$, $\cosh$, $\tanh$, $\coth$}

\subsection*{Notation and preliminaries}
Throughout, write $z=x+iy$, $x,y\in\mathbb R$, and $w=u+iv$ for the image.
Recall the basic identities
\[
\cosh z=\frac{e^{z}+e^{-z}}{2},\qquad
\sinh z=\frac{e^{z}-e^{-z}}{2},\qquad
\tanh z=\frac{\sinh z}{\cosh z},\qquad
\coth z=\frac{\cosh z}{\sinh z}.
\]
We use the principal branch of the logarithm $\log$ on $\mathbb C\setminus(-\infty,0]$, and principal square roots $\sqrt{\cdot}$ with branch cut on $(-\infty,0]$ unless stated otherwise.

\bigskip\hrule\bigskip

% ============================ 1. SINH ============================
\subsection*{1.\ $w=\sinh z$}

\paragraph{Analytic background.}\mbox{}\\[2pt]
$\sinh$ is an \emph{entire}, \emph{odd} function, with imaginary-periodicity:
\[
\sinh(-z)=-\sinh z,\qquad \sinh(z+2\pi i)=\sinh z.
\]
Zeros occur at $z_k=i\pi k$ for $k\in\mathbb Z$ (all simple). No poles.

\medskip

\paragraph{Cauchy--Riemann decomposition.}\mbox{}\\[2pt]
With $z=x+iy$,
\[
\sinh z=\sinh x\cos y\ +\ i\,\cosh x\sin y \quad\Rightarrow\quad
u=\sinh x\cos y,\ \ v=\cosh x\sin y.
\]
From this, one obtains orthogonal families of confocal conics:
\begin{align*}
	x=\text{const}&:\quad \left(\frac{u}{\sinh x}\right)^{2}+\left(\frac{v}{\cosh x}\right)^{2}=1
	&&\text{(ellipses, foci at $\pm i$)},\\[2pt]
	y=\text{const}&:\quad \left(\frac{u}{\cos y}\right)^{2}-\left(\frac{v}{\sin y}\right)^{2}=-1
	&&\text{(hyperbolas, foci at $\pm i$)}.
\end{align*}
Real and imaginary axes map as:
\[
y=0:\ w=\sinh x\in\mathbb R\ \text{(onto $\mathbb R$)},\qquad
x=0:\ w=i\sin y\ \text{(onto the segment $i[-1,1]$)}.
\]

\medskip

\paragraph{Jacobian and conformality.}\mbox{}\\[2pt]
$\sinh'(z)=\cosh z$, and $\cosh z\neq 0$ on the strip $|\Im z|<\tfrac{\pi}{2}$; hence $\sinh$ is conformal there.

\medskip

\paragraph{Canonical strip mapping (global picture).}\mbox{}\\[2pt]
\[
\boxed{\
	\sinh:\ \{\,|\Im z|<\tfrac{\pi}{2}\,\}\ \xrightarrow[\ \text{biholomorphism}\ ]{}\ 
	\mathbb C\setminus i\big((-\infty,-1]\cup[1,\infty)\big)\ .
}
\]
\emph{Explanation.} The boundary lines $y=\pm\frac{\pi}{2}$ map to the imaginary rays $\pm i[1,\infty)$
(via $v=\cosh x\,\sin y$ with $\sin(\pm\pi/2)=\pm1$ and $u=\sinh x\cos(\pm\pi/2)=0$), while the interior maps injectively since $\cosh z\neq0$ and periodicity in the imaginary direction has period $2\pi$, not $\pi$.

\medskip

\paragraph{Inverse branch and branch cuts.}\mbox{}\\[2pt]
A principal inverse on the slit domain is
\[
\operatorname{arsinh} w=\log\!\big(w+\sqrt{w^{2}+1}\big),
\]
with branch cuts along $i(-\infty,-1]\cup i[1,\infty)$ so that $\Im(\operatorname{arsinh} w)\in(-\tfrac{\pi}{2},\tfrac{\pi}{2})$.

\medskip

\paragraph{Useful identities.}\mbox{}\\[2pt]
\[
\cosh^{2}z-\sinh^{2}z=1,\qquad
\sinh(z+\zeta)=\sinh z\,\cosh\zeta+\cosh z\,\sinh\zeta.
\]

\bigskip\hrule\bigskip

% ============================ 2. COSH ============================
\subsection*{2.\ $w=\cosh z$}

\paragraph{Analytic background.}\mbox{}\\[2pt]
$\cosh$ is \emph{entire}, \emph{even}, and $2\pi i$-periodic:
\[
\cosh(-z)=\cosh z,\qquad \cosh(z+2\pi i)=\cosh z.
\]
Zeros at $z=i\pi\big(k+\tfrac12\big)$, all simple. No poles.

\medskip

\paragraph{Cauchy--Riemann decomposition.}\mbox{}\\[2pt]
With $z=x+iy$,
\[
\cosh z=\cosh x\cos y\ +\ i\,\sinh x\sin y \quad\Rightarrow\quad
u=\cosh x\cos y,\ \ v=\sinh x\sin y.
\]
Again one gets confocal conics:
\begin{align*}
	x=\text{const}&:\quad \left(\frac{u}{\cosh x}\right)^{2}+\left(\frac{v}{\sinh x}\right)^{2}=1
	&&\text{(ellipses, foci at $\pm1$)},\\[2pt]
	y=\text{const}&:\quad \left(\frac{u}{\cos y}\right)^{2}-\left(\frac{v}{\sin y}\right)^{2}=1
	&&\text{(hyperbolas, foci at $\pm1$)}.
\end{align*}
On the axes,
\[
y=0:\ w=\cosh x\in[1,\infty),\qquad
x=0:\ w=\cos y\in[-1,1].
\]

\medskip

\paragraph{Critical points and injectivity on strips.}\mbox{}\\[2pt]
$\cosh'(z)=\sinh z$, which vanishes at $z=i\pi k$. Thus $\cosh$ fails to be injective across lines $y=k\pi$, but is injective on any strip avoiding these lines (e.g.\ $0<\Im z<\pi$).

\medskip

\paragraph{Canonical strip mapping.}\mbox{}\\[2pt]
\[
\boxed{\
	\cosh:\ \{\,0<\Im z<\pi\,\}\ \xrightarrow[\ \text{biholomorphism}\ ]{}\ 
	\mathbb C\setminus\big((-\infty,-1]\cup[1,\infty)\big)\ .
}
\]
\emph{Explanation.} The midline $y=0$ maps to $[1,\infty)$; the top boundary $y=\pi$ maps to $(-\infty,-1]$; injectivity holds on $0<\Im z<\pi$ since there $\sinh z\neq0$.

\medskip

\paragraph{Inverse branch and branch cuts.}\mbox{}\\[2pt]
\[
\operatorname{arcosh} w=\log\!\Big(w+\sqrt{w-1}\,\sqrt{w+1}\Big),
\]
with branch cuts along $(-\infty,-1]\cup[1,\infty)$ so that $\Re(\operatorname{arcosh} w)\ge0$ and $0<\Im(\operatorname{arcosh} w)<\pi$ on the principal sheet.

\medskip

\paragraph{Useful identities.}\mbox{}\\[2pt]
\[
\cosh(z+\zeta)=\cosh z\,\cosh\zeta+\sinh z\,\sinh\zeta,\qquad
\cosh z=1+2\sinh^{2}\!\frac{z}{2}.
\]

\bigskip\hrule\bigskip

% ============================ 3. TANH ============================
\subsection*{3.\ $w=\tanh z$}

\paragraph{Analytic background.}\mbox{}\\[2pt]
$\tanh$ is \emph{meromorphic}, \emph{odd}, and $i\pi$-periodic:
\[
\tanh(-z)=-\tanh z,\qquad \tanh(z+i\pi)=\tanh z.
\]
Zeros at $z=i\pi k$; poles (simple) at $z=i\pi\big(k+\tfrac12\big)$.

\medskip

\paragraph{Strip $\leftrightarrow$ disk (Cayley-type).}\mbox{}\\[2pt]
\[
\boxed{\
	\tanh:\ \{\,|\Im z|<\tfrac{\pi}{2}\,\}\ \xrightarrow[\ \text{biholomorphism}\ ]{}\ \mathbb D\ .
}
\]
\emph{Boundary correspondence.} The real axis $y=0$ maps bijectively onto $(-1,1)$ (strictly increasing).
As $y\to\pm\frac{\pi}{2}$, $|\tanh(x+iy)|\to1$ uniformly in $x$; hence the horizontal boundaries map to $\partial\mathbb D$.

\medskip

\paragraph{Inverse (disk $\to$ strip).}\mbox{}\\[2pt]
\[
\operatorname{artanh} w=\frac{1}{2}\log\!\left(\frac{1+w}{1-w}\right),
\qquad |w|<1,
\]
which maps $\mathbb D$ conformally onto $\{|\Im z|<\tfrac{\pi}{2}\}$.

\medskip

\paragraph{Derivative and Schwarz--Pick.}\mbox{}\\[2pt]
\[
(\tanh z)'=\operatorname{sech}^{2}\!z,\qquad
\left|\tanh'(z)\right|=\frac{1}{|\cosh z|^{2}}
\]
and conjugating by a disk automorphism shows $\tanh$ is extremal for the disk-strip version of Schwarz--Pick.

\medskip

\paragraph{Level sets.}\mbox{}\\[2pt]
Lines $x=\text{const}$ and $y=\text{const}$ in the strip map to orthogonal circle arcs inside $\mathbb D$ meeting $\partial\mathbb D$ orthogonally (images of preimages under the conformal equivalence with the half-plane followed by a Cayley map).

\bigskip\hrule\bigskip

% ============================ 4. COTH ============================
\subsection*{4.\ $w=\coth z$}

\paragraph{Analytic background.}\mbox{}\\[2pt]
$\coth$ is \emph{meromorphic}, \emph{odd}, and $i\pi$-periodic:
\[
\coth(-z)=-\coth z,\qquad \coth(z+i\pi)=\coth z.
\]
Poles (simple) at $z=i\pi k$; zeros (simple) at $z=i\pi\big(k+\tfrac12\big)$.

\medskip

\paragraph{Strip $\leftrightarrow$ exterior disk.}\mbox{}\\[2pt]
\[
\boxed{\
	\coth:\ \{\,|\Im z|<\tfrac{\pi}{2}\,\}\setminus\{0\}\ \xrightarrow[\ \text{biholomorphism}\ ]{}\ \{\,|w|>1\,\}\ .
}
\]
\emph{Boundary correspondence.} The real axis (excluding $0$) maps to the rays $(-\infty,-1)\cup(1,\infty)$.
As $y\to\pm\frac{\pi}{2}$, $|\coth(x+iy)|\to1^{+}$; thus horizontal boundaries map to the unit circle from the \emph{exterior} side.

\medskip

\paragraph{Inverse (exterior disk $\to$ strip).}\mbox{}\\[2pt]
\[
\operatorname{arcoth} w=\frac{1}{2}\log\!\left(\frac{w+1}{w-1}\right),\qquad |w|>1,
\]
with the standard logarithm branch so that $\Im(\operatorname{arcoth} w)\in(-\tfrac{\pi}{2},\tfrac{\pi}{2})$.

\medskip

\paragraph{Relation with $\tanh$.}\mbox{}\\[2pt]
\[
\coth z=\tanh\!\left(z+\tfrac{i\pi}{2}\right),
\]
so the exterior-disk picture is just the disk picture shifted vertically by $i\pi/2$ and inverted via $w\mapsto 1/w$.

\bigskip\hrule\bigskip

% ============================ 5. Global relations & composition ============================
\subsection*{5.\ Global relations, inverse formulas, and composition patterns}

\paragraph{Inverse formulas (principal branches).}\mbox{}\\[2pt]
\[
\operatorname{arsinh} w=\log\!\big(w+\sqrt{w^{2}+1}\big),\qquad
\operatorname{arcosh} w=\log\!\big(w+\sqrt{w-1}\,\sqrt{w+1}\big),
\]
\[
\operatorname{artanh} w=\frac12\log\!\left(\frac{1+w}{1-w}\right),\qquad
\operatorname{arcoth} w=\frac12\log\!\left(\frac{w+1}{w-1}\right).
\]
Branch cuts are chosen to match the canonical target domains described above.

\medskip

\paragraph{Disk/half-plane equivalences via Cayley.}\mbox{}\\[2pt]
Using the Cayley map $C(\zeta)=\frac{1+\zeta}{1-\zeta}$ and its inverse $C^{-1}(w)=\frac{w-1}{w+1}$,
\[
\tanh z=\ C^{-1}\!\left(e^{2z}\right),\qquad
\coth z=\ C^{-1}\!\left(e^{2z}\right)^{-1},\qquad
\sinh z=\frac{e^{z}-e^{-z}}{2},\quad \cosh z=\frac{e^{z}+e^{-z}}{2}.
\]
Thus strip $\leftrightarrow$ disk/exterior-disk mappings for $\tanh$ and $\coth$ are just the exponential’s strip-to-sector equivalences composed with a Cayley map.

\medskip

\paragraph{Growth and boundary behavior.}\mbox{}\\[2pt]
As $x\to\pm\infty$ with $y$ fixed, $\sinh z\sim \frac12 e^{x}e^{iy}$ and $\cosh z\sim \frac12 e^{x}e^{iy}$; hence horizontal lines map to asymptotically radial curves. For $\tanh$ and $\coth$, $|\tanh(x+iy)|\to 1$ and $|\coth(x+iy)|\to 1$ exponentially in $|x|$ (uniform in $y$ within the strip height).

\medskip

\paragraph{Critical sets and injectivity regions.}\mbox{}\\[2pt]
$\sinh'(z)=\cosh z$ vanishes along $y=\frac{\pi}{2}+\pi\mathbb Z$; $\cosh'(z)=\sinh z$ vanishes along $y=\pi\mathbb Z$.
Therefore, the natural injectivity strips are $|\Im z|<\tfrac{\pi}{2}$ for $\sinh$ and $0<\Im z<\pi$ for $\cosh$.
For $\tanh$ and $\coth$, poles on the horizontal boundaries determine the natural strip height $\pi$.

\bigskip\hrule\bigskip

% ============================ 6. Summary boxes ============================
\subsection*{6.\ Summary: canonical conformal equivalences (boxed)}
\[
\boxed{\
	\sinh:\ \{\,|\Im z|<\tfrac{\pi}{2}\,\}\ \longrightarrow\ 
	\mathbb C\setminus i\big((-\infty,-1]\cup[1,\infty)\big)\ \ \text{(biholomorphism)}\ .
}
\]
\[
\boxed{\
	\cosh:\ \{\,0<\Im z<\pi\,\}\ \longrightarrow\ 
	\mathbb C\setminus\big((-\infty,-1]\cup[1,\infty)\big)\ \ \text{(biholomorphism)}\ .
}
\]
\[
\boxed{\
	\tanh:\ \{\,|\Im z|<\tfrac{\pi}{2}\,\}\ \longrightarrow\ \mathbb D\ \ \text{(biholomorphism)}\ .
}
\]
\[
\boxed{\
	\coth:\ \{\,|\Im z|<\tfrac{\pi}{2}\,\}\setminus\{0\}\ \longrightarrow\ \{\,|w|>1\,\}\ \ \text{(biholomorphism)}\ .
}
\]

\medskip

\subsection*{7.\ (Optional) Proof sketches for strip mappings}

\paragraph{Injectivity for $\sinh$ on $|\Im z|<\pi/2$.}
If $\sinh z_{1}=\sinh z_{2}$ then $e^{z_{1}}-e^{-z_{1}}=e^{z_{2}}-e^{-z_{2}}$.
Rearranging gives $\{e^{z_{1}},e^{-z_{1}}\}=\{e^{z_{2}},e^{-z_{2}}\}$.
Taking logs in the horizontal strip $|\Im z|<\pi/2$ (where the exponential is injective modulo $2\pi i$) forces $z_{1}=z_{2}$.

\paragraph{Boundary images for $\sinh$.}
Set $z=x\pm i\pi/2$. Then $\sinh z=\sinh x\cdot 0 \ \pm i\,\cosh x$, so the boundaries map to $\pm i[1,\infty)$; the interior cannot reach those rays by the open mapping theorem and the maximum modulus principle applied to $1/(\sinh z\mp i)$.

\paragraph{Analogous arguments} apply to $\cosh$ (use $\sinh$ zeros to avoid critical lines) and to $\tanh,\coth$ using the representation via exponentials and the Cayley transform.




% =====================================================================
% Inverse Circular Functions: arcsin, arccos, arctan, arccot
% Principal branches, branch cuts, strip images, boundary correspondence
% =====================================================================

\section*{Inverse Circular Functions: $\arcsin$, $\arccos$, $\arctan$, $\arccot$}

\subsection*{Conventions}
We use principal branches of $\log$ and $\sqrt{\cdot}$ with branch cut $(-\infty,0]$ unless stated otherwise.
For each inverse, the principal domain is obtained by slitting the $w$–plane along the images of the multiple values.

\bigskip\hrule\bigskip

% ============================ arcsin ============================
\subsection*{1.\ $z=\arcsin w$}

\paragraph{Principal formula (single-valued on the slit plane).}\mbox{}\\[2pt]
\BoxEq{ \arcsin w = -\,i\,\log\!\Big(i\,w+\sqrt{\,1-w^{2}\,}\Big) }

\paragraph{Branch points \& cuts.}\mbox{}\\[2pt]
Branch points at $w=\pm1$; principal cuts along $(-\infty,-1]\cup[1,\infty)$.

\paragraph{Target strip (canonical range).}\mbox{}\\[2pt]
\BoxEq{ \arcsin:\ \mathbb C\setminus\big((-\infty,-1]\cup[1,\infty)\big)\ \longrightarrow\ \{\,|\Re z|<\tfrac{\pi}{2}\,\} }
For $w\in(-1,1)$, $\arcsin w\in(-\tfrac{\pi}{2},\tfrac{\pi}{2})\subset\mathbb R$.
For $w>1$,
$\arcsin w=\tfrac{\pi}{2}-i\,\operatorname{arcosh} w$ (hence $\Re=\tfrac{\pi}{2}$);
for $w<-1$,
$\arcsin w=-\tfrac{\pi}{2}+i\,\operatorname{arcosh}(-w)$ (hence $\Re=-\tfrac{\pi}{2}$).

\paragraph{Boundary correspondence.}\mbox{}\\[2pt]
Approaching the cuts $[1,\infty)$ or $(-\infty,-1]$ forces $\Re \arcsin w\to \pm\tfrac{\pi}{2}$ and $|\Im \arcsin w|\to\infty$.
The real axis segment $(-1,1)$ maps to the real segment $(-\tfrac{\pi}{2},\tfrac{\pi}{2})$.

\bigskip\hrule\bigskip

% ============================ arccos ============================
\subsection*{2.\ $z=\arccos w$}

\paragraph{Principal formulae.}\mbox{}\\[2pt]
\BoxEq{ \arccos w = \frac{\pi}{2}-\arcsin w }
\BoxEq{ \arccos w = \frac{1}{i}\,\log\!\Big(w+\sqrt{w-1}\,\sqrt{w+1}\Big) }

\paragraph{Branch points \& cuts.}\mbox{}\\[2pt]
Same as $\arcsin$: branch points $\pm1$; cuts $(-\infty,-1]\cup[1,\infty)$.

\paragraph{Target strip (canonical range).}\mbox{}\\[2pt]
\BoxEq{ \arccos:\ \mathbb C\setminus\big((-\infty,-1]\cup[1,\infty)\big)\ \longrightarrow\ \{\,0<\Re z<\pi\,\} }
For $w\in(-1,1)$, $\arccos w\in(0,\pi)\subset\mathbb R$.
For $w>1$, $\arccos w=i\,\operatorname{arcosh} w$ (so $\Re=0$);
for $w<-1$, $\arccos w=\pi - i\,\operatorname{arcosh}(-w)$ (so $\Re=\pi$).

\paragraph{Boundary correspondence.}\mbox{}\\[2pt]
The two slits map to the vertical strip boundaries $\Re z=0$ and $\Re z=\pi$.
The central interval $(-1,1)$ maps to the real segment $(0,\pi)$.

\bigskip\hrule\bigskip

% ============================ arctan ============================
\subsection*{3.\ $z=\arctan w$}

\paragraph{Principal formula.}\mbox{}\\[2pt]
\BoxEq{ \arctan w = \frac{1}{2i}\,\log\!\left(\frac{1+i\,w}{\,1-i\,w\,}\right) }

\paragraph{Branch points \& cuts.}\mbox{}\\[2pt]
Branch points at $w=\pm i$; principal cuts $i(-\infty,-1]\cup i[1,\infty)$ along the imaginary axis beyond $\pm i$.

\paragraph{Target strip (canonical range).}\mbox{}\\[2pt]
\BoxEq{ \arctan:\ \mathbb C\setminus i\big((-\infty,-1]\cup[1,\infty)\big)\ \longrightarrow\ \{\,|\Re z|<\tfrac{\pi}{2}\,\} }
On $\mathbb R$, $\arctan \mathbb R=(-\tfrac{\pi}{2},\tfrac{\pi}{2})$.
Approaching $w=\pm i\,t$ ($t\ge1$) drives $\Re \arctan w\to \pm\tfrac{\pi}{2}$ and $|\Im|\to\infty$.

\paragraph{Symmetries.}\mbox{}\\[2pt]
$\arctan(-w)=-\arctan w$; conjugation across the real axis preserves values away from cuts.

\bigskip\hrule\bigskip

% ============================ arccot ============================
\subsection*{4.\ $z=\arccot w$}

\paragraph{Principal formulae.}\mbox{}\\[2pt]
\BoxEq{ \arccot w = \frac{1}{2i}\,\log\!\left(\frac{w+i}{\,w-i\,}\right) }
\BoxEq{ \arccot w = \frac{\pi}{2}-\arctan w }

\paragraph{Branch points \& cuts.}\mbox{}\\[2pt]
As for $\arctan$: branch points at $w=\pm i$; cuts $i(-\infty,-1]\cup i[1,\infty)$.

\paragraph{Target strip (canonical range).}\mbox{}\\[2pt]
\BoxEq{ \arccot:\ \mathbb C\setminus i\big((-\infty,-1]\cup[1,\infty)\big)\ \longrightarrow\ \{\,0<\Re z<\pi\,\} }
On $\mathbb R$, $\arccot \mathbb R=(0,\pi)$.
Approaching the cuts pins the image to $\Re z=0$ or $\Re z=\pi$ with $|\Im|\to\infty$.

\bigskip\hrule\bigskip

% ============================ Relations & identities ============================
\subsection*{5.\ Relations and identities (principal branches)}
\[
\arcsin w + \arccos w = \frac{\pi}{2},\qquad
\arctan w + \arccot w = \frac{\pi}{2}.
\]
\[
\arcsin w = -\,i\,\operatorname{arsinh}(i w),\qquad
\arccos w = \operatorname{arcosh} w \cdot \frac{1}{i}\ \ (\text{on }[1,\infty)).
\]
\[
\arctan w = \frac{1}{2i}\big(\log(1+i w)-\log(1-i w)\big),\quad
\arccot w = \frac{1}{2i}\big(\log(w+i)-\log(w-i)\big).
\]

\medskip

\subsection*{6.\ Boundary/axis behavior (summary)}
\begin{itemize}
	\item $\arcsin$: real on $(-1,1)$; the cuts map to $\Re z=\pm\tfrac{\pi}{2}$.
	\item $\arccos$: real on $(-1,1)$; the cuts map to $\Re z=0,\ \pi$.
	\item $\arctan$: real on $\mathbb R$; the imaginary-axis cuts beyond $\pm i$ map to $\Re z=\pm\tfrac{\pi}{2}$.
	\item $\arccot$: real on $\mathbb R$; the same cuts map to $\Re z=0,\ \pi$.
\end{itemize}






	
\section*{Conformal maps: an overview catalogue (no details)}

\begin{enumerate}
	\item \textbf{Cayley (disk $\leftrightarrow$ right half-plane).}\\[2pt]
	\[
	\boxed{\,w=\frac{1+z}{1-z}\,}\qquad\text{inverse:}\quad z=\frac{w-1}{w+1}.
	\]
	
	\item \textbf{Half-plane $\to$ half-plane (three-point normalization).}\\[2pt]
	\[
	\boxed{\,w=\frac{(z-a)(c-b)}{(z-b)(c-a)}\,}.
	\]
	
	\item \textbf{Disk automorphisms.}\\[2pt]
	\[
	\boxed{\,\phi_{a,\theta}(z)=e^{i\theta}\frac{z-a}{1-\bar a z}\,},\qquad |a|<1.
	\]
	
	\item \textbf{Strip $|\Im z|<h/2$ $\to$ right half-plane.}\\[2pt]
	\[
	\boxed{\,w=\exp\!\left(\frac{\pi z}{h}\right)\,}.
	\]
	
	\item \textbf{Strip $|\Im z|<h/2$ $\to$ unit disk.}\\[2pt]
	\[
	\boxed{\,\zeta=\frac{e^{\pi z/h}-1}{\,e^{\pi z/h}+1\,}\,}.
	\]
	
	\item \textbf{Strip $\to$ strip (height scaling).}\\[2pt]
	\[
	\boxed{\,z\mapsto \frac{h_2}{h_1}\,z\,}.
	\]
	
	\item \textbf{Band $a<\Im z<b$ $\to$ centered strip.}\\[2pt]
	\[
	\boxed{\,z\mapsto \frac{\pi}{b-a}\!\left(z-i\frac{a+b}{2}\right)\,}.
	\]
	
	\item \textbf{Sector $\{0<\arg z<\alpha\}$ $\to$ upper half-plane.}\\[2pt]
	\[
	\boxed{\,w=z^{\pi/\alpha}\,}\ \ \text{(principal branch)}.
	\]
	
	\item \textbf{Wedge $\to$ disk.}\\[2pt]
	Compose item (8) with Cayley (item 1).
	
	\item \textbf{Upper half-plane $\to$ disk (prescribed triple).}\\[2pt]
	\[
	\boxed{\,\Phi(z)=\frac{(z-a)}{(z-\bar a)}\cdot \frac{(b-\bar a)}{(b-a)}\,}.
	\]
	
	\item \textbf{Annulus $r<|z|<1$ $\leftrightarrow$ strip.}\\[2pt]
	\[
	\boxed{\,w=\log z\,}\quad\text{maps to }\ \log r<\Re w<0.
	\]
	
	\item \textbf{Annulus $\to$ disk.}\\[2pt]
	Use $w=\log z$ then item (5).
	
	\item \textbf{Half-plane minus disk $\to$ annulus.}\\[2pt]
	\[
	\boxed{\,\zeta=\frac{z-\alpha}{z+\alpha}\,},\qquad \alpha^{2}=a^{2}-b^{2}.
	\]
	
	\item \textbf{Disk with off-center hole $\to$ annulus.}\\[2pt]
	Center via $\phi_a$ (item 3), then items (11)–(12).
	
	\item \textbf{Punctured plane $\leftrightarrow$ cylinder.}\\[2pt]
	\[
	\boxed{\,w=\log z\,}\quad(\text{period }2\pi i).
	\]
	
	\item \textbf{Strip $\to$ punctured plane.}\\[2pt]
	\[
	\boxed{\,z\mapsto e^{z}\,}.
	\]
	
	\item \textbf{Ray-slit plane $\mathbb C\setminus[0,\infty)$ $\to$ half-plane.}\\[2pt]
	\[
	\boxed{\,w=\sqrt{z}\,}\ \ \text{(choose branch)}.
	\]
	
	\item \textbf{Segment-slit plane $\mathbb C\setminus[a,b]$ $\to$ half-plane.}\\[2pt]
	\[
	\boxed{\,w=\sqrt{\frac{z-a}{\,z-b\,}}\,}.
	\]
	
	\item \textbf{Exterior of $[-1,1]$ $\leftrightarrow$ exterior unit circle (Joukowski).}\\[2pt]
	\[
	\boxed{\,w=z+\frac{1}{z}\,},\qquad
	z=\frac{w\pm\sqrt{w^{2}-4}}{2}.
	\]
	
	\item \textbf{Exterior of ellipse $\leftrightarrow$ exterior unit circle.}\\[2pt]
	\[
	\boxed{\,w=\tfrac12\!\left(c z+\frac{1}{c z}\right)\,}\quad(\text{$c$ from axes}).
	\]
	
	\item \textbf{Right half-plane $\to$ vertical strip.}\\[2pt]
	\[
	\boxed{\,z\mapsto \frac{1}{\pi}\log\frac{z-1}{\,z+1\,}\,}.
	\]
	
	\item \textbf{Upper half-plane $\to$ rectangle (SC/elliptic).}\\[2pt]
	\[
	\boxed{\,\zeta=\int^z \frac{dt}{\sqrt{(t-a)(t-b)(t-c)(t-d)}}\,}.
	\]
	
	\item \textbf{Rectangle $\to$ disk (Jacobi sn).}\\[2pt]
	\[
	\boxed{\,\zeta=\operatorname{sn}\!\left(\frac{K(k)}{L}\,z;\,k\right)\,}.
	\]
	
	\item \textbf{Half-disk $\leftrightarrow$ half-plane.}\\[2pt]
	Restriction of Cayley (item 1) and inverse.
	
	\item \textbf{Lens (intersection of two disks) $\to$ disk.}\\[2pt]
	Möbius normalizes the circles; then a power map; then Cayley.
	
	\item \textbf{Half-strip $\{x>0,\ 0<\Im z<h\}$ $\to$ half-disk.}\\[2pt]
	\[
	\boxed{\,\exp(\pi z/h)\,}\ \ \text{then power, then Cayley}.
	\]
	
	\item \textbf{Polygon $\to$ half-plane (Schwarz--Christoffel).}\\[2pt]
	\[
	\boxed{\,\Phi'(z)=C\prod_k (z-z_k)^{\alpha_k-1}\,},\qquad \sum\alpha_k=2.
	\]
	
	\item \textbf{Upper half-plane $\to$ slit half-plane.}\\[2pt]
	\[
	\boxed{\,w=z+\frac{1}{z}\,}\ \ \text{then a real Möbius adjustment}.
	\]
	
	\item \textbf{Disk $\to$ disk with boundary arc removed.}\\[2pt]
	Automorphism to put endpoints at $\pm1$, then boundary power $z^\lambda$.
	
	\item \textbf{Finite Blaschke product (disk $\to$ disk, degree $n$).}\\[2pt]
	\[
	\boxed{\,B(z)=e^{i\theta}\prod_{k=1}^n \frac{z-a_k}{\,1-\bar a_k z\,}\,}.
	\]
\end{enumerate}

	
	
	
\section*{1.\ Cayley transform: $\mathbb D \leftrightarrow$ right half-plane}

\paragraph{Map and inverse.}\mbox{}\\[2pt]
\[
w=\frac{1+z}{\,1-z\,}
\quad\text{maps}\quad
\mathbb D=\{|z|<1\}\ \text{onto}\ \{\,\Re w>0\,\},
\qquad
z=\frac{w-1}{\,w+1\,}
\]
is the inverse.

\medskip

\paragraph{Why $\Re w>0$.}\mbox{}\\[2pt]
Compute
\[
\Re\!\left(\frac{1+z}{\,1-z\,}\right)
=\frac{\,1-|z|^{2}\,}{\,|1-z|^{2}\,}.
\]
Thus $\Re w>0$ for $|z|<1$, and $\Re w=0$ for $|z|=1$ (except $z=1$, which maps to $w=\infty$).
Therefore
\[
|z|<1\ \Longleftrightarrow\ \Re w>0,
\qquad
|z|=1\ \Longleftrightarrow\ \Re w=0 .
\]

\medskip

\paragraph{Boundary correspondence.}\mbox{}\\[2pt]
$z=0\mapsto w=1$,
\quad
$z\to1^{-}\mapsto w\to+\infty$,
\quad
and $\partial\mathbb D\setminus\{1\}$ maps to the line $\Re w=0$.\\[2pt]
Circles/lines orthogonal to $\partial\mathbb D$ map to \emph{vertical} lines in the $w$–plane.

\bigskip\hrule\bigskip

\section*{2.\ Three-point Möbius normalization: half-plane $\leftrightarrow$ half-plane}

\paragraph{Cross-ratio map to the real axis.}\mbox{}\\[2pt]
Let $L$ be a line or circle, and choose three distinct points $a,b,c\in L$. Define
\[
\Phi(z)=\frac{(z-a)(c-b)}{(z-b)(c-a)}.
\]
Then $\Phi(L)\subset\mathbb R\cup\{\infty\}$, with
\[
\Phi(a)=0,\qquad \Phi(b)=\infty,\qquad \Phi(c)=1,
\]
and the two sides of $L$ map to the half-planes $\Im\Phi(z)\gtrless 0$.

\medskip

\paragraph{Half-plane to half-plane via three boundary points.}\mbox{}\\[2pt]
Given half-planes $H_1,H_2$ with boundary arcs $L_1,L_2$,
pick ordered boundary triples $a,b,c\in L_1$ and $A,B,C\in L_2$ (same cyclic order).
Then
\[
T(z)=\frac{(z-a)(c-b)}{(z-b)(c-a)}\cdot\frac{(C-A)}{(C-B)}
\]
is a Möbius map sending $L_1\!\to L_2$ with
\[
a\mapsto A,\qquad b\mapsto B,\qquad c\mapsto C,
\]
and taking the chosen side of $H_1$ onto the chosen side of $H_2$.

\medskip

\paragraph{Specialization to the real line.}\mbox{}\\[2pt]
To send $L$ to $\mathbb R$, use $\Phi$ above; then pre/postcompose with an \emph{affine real} map
\[
\xi\longmapsto \alpha\,\xi+\beta,\qquad \alpha>0,
\]
to obtain any target half-plane, e.g.\ $\Re w>0$ or $\Im w>0$.

	
	
\end{document}

